\subsection{Top-down control and neuromodulation}

In the biological circuit, Si3 neurons are not intrinsically oscillatory but receive tonic excitatory drive from Si1 command neurons that initiate and sustain swim episodes \cite{sakurai2019command}.
We now examine a five-cell network configuration that incorporates Si1 control, with Si3 neurons set to quiescent baseline states that require Si1 excitation to become active.

\subsubsection{Si1 gating of network activity}

In \textit{Dendronotus iris}, Si1 neurons function as command-like neurons that initiate and terminate swim episodes \cite{sakurai2019command}.
They provide tonic excitatory drive to both Si3 neurons, which in turn activate the four-cell CPG through the contralateral E/I modules.
This hierarchical organization allows Si1 neurons to gate network state transitions while intrinsic CPG dynamics shape burst structure and left-right alternation.

Figure~\ref{fig:si1_neuromod}A summarizes the network architecture used for this analysis.
Si1 provides excitatory drive to both Si3 neurons and also modulates the Si3$\to$Si2 excitatory pathway, thereby coupling episode recruitment to burst-frequency regulation within the same command system.

\subsubsection{Neuromodulatory burst frequency control}

Beyond gating network activity, Si1 neurons regulate burst frequency in an activity-dependent manner.
Experimental recordings demonstrate that swim frequency increases with Si1 firing rate, establishing that the command signal encodes not merely on/off state but graded frequency information \cite{sakurai2019command}.
This observation motivates our exploration of how neuromodulation of synaptic properties in the Si3$\to$Si2 excitatory pathway could implement such graded control.
We compared two distinct mechanisms: presynaptic modulation of the logistic Si3$\to$Si2 accumulation-rate parameter $\alpha$, and postsynaptic modulation of the maximal effective conductance in the same pathway.

Figure~\ref{fig:si1_neuromod}B compares the biological Si1-rate versus burst-rate relationship with matched measurements extracted from representative model traces.
Both model variants capture the positive relationship between command-neuron firing and burst frequency, but the presynaptic model occupies the biological range more naturally across successive burst intervals, whereas the postsynaptic model is compressed toward higher burst frequencies for a broad range of Si1 firing rates.

Figure~\ref{fig:si1_neuromod}C shows the full gain scan during ongoing drive.
Presynaptic modulation produces a broader and smoother increase in burst frequency as gain increases.
Postsynaptic modulation also increases burst frequency, but the response is flatter over substantial parts of the scan and reaches a comparatively narrow high-frequency band.
Thus, both mechanisms can tune the active rhythm, but presynaptic modulation provides a more graded control landscape.

Figure~\ref{fig:si1_neuromod}E shows representative biological and model traces.
The biological recording establishes the empirical relationship between elevated Si1 firing and Si2 bursting, while the two model traces illustrate how the same Si1 drive can be translated into different burst statistics depending on whether modulation acts on release kinetics or effective conductance.

Taken together, these results support a restrained conclusion: Si1-mediated modulation of the Si3$\to$Si2 pathway is sufficient to regulate burst frequency during active swimming, and modulation of synaptic kinetics provides the closer match to the observed graded Si1-to-burst relationship.
This generates a testable prediction that pharmacological manipulations targeting presynaptic release machinery should produce smoother frequency tuning than manipulations acting primarily on postsynaptic conductance.
